% =====================================================================================================
\section{Related Work}\label{sec:related}
\subsection{Selective encryption of JPEG}
Van Droogenbroeck and Benedett encrypted the appended bits of selected AC coefficients of baseline JPEG, leaving the Huffman code words in place, so that the stream stays compliant at a constant bit rate~\cite{vandroogenbroeck2002}. Puech and Rodrigues encrypted selected DCT coefficients of medical images with AES in OFB mode~\cite{puech2005}. Auer et al.\ encrypted baseline JPEG at bitstream level, permuting Huffman code-word/value pairs and toggling value bits, with an AES-OFB generator at more than 25 VGA frames per second, preserving length and format~\cite{auer2013}. Massoudi et al.\ survey the field and note that visual distortion is subjective and does not measure security: replacement and chosen-plaintext attacks recover more than the scrambled image suggests~\cite{massoudi2008}. Selective encryption has also been extended to H.264/AVC and HEVC~\cite{boyadjis2017}.

\subsection{Attacks on format-compliant encryption}
Lin et al.\ show that DCT sign bits can be recovered by linear programming~\cite{lin2024}, which bounds what sign-only encryption can promise. Hirose et al.\ analyse bitstream-level JPEG encryption that shuffles blocks delimited by restart markers; the number of non-zero coefficients per block, which no field encryption hides, recovers contours unless the segments are permuted, and partly again when the restart interval is large~\cite{hirose2024}. Our strongest level uses the same mechanism with a keyed, per-frame permutation. Li et al.\ recover block-shuffled JPEG with a designed plaintext that reveals the permutation through its per-block non-zero counts, which then inverts any image shuffled with the same permutation~\cite{li2025shuffle}. Learned attackers are a threat beyond hand-made ones: Sirichotedumrong and Kiya find that a GAN-based ciphertext-only attack recovers visual content from most image-protection schemes designed for privacy-preserving deep learning~\cite{sirichotedumrong2020}. Block permutations also invite jigsaw solvers, from greedy placement~\cite{pomeranz2011} to minimum-spanning-tree assembly with gradient compatibility~\cite{gallagher2012}.

\subsection{Hardware JPEG encoders with encryption}
Kishore et al.\ added PN-sequence block and symbol scrambling to an FPGA JPEG encoder for 256$\times$256 grey images on Stratix~II and~IV; they describe the extra hardware as about 1\,\%, while their tables show 13--17\,\% more ALUTs, twice the memory bits and 9--10\,\% fewer frames per second with it~\cite{kishore2015}. Ba\u{g}baba and \"Ors combined an FPGA JPEG encoder with TEA encryption of the DC coefficients, or of the DC and first five AC coefficients, of every block~\cite{bagbaba2015}. Neither authenticates the frames, and both use key material that does not change from frame to frame. Qiu et al.\ XOR raw video frames with two one-dimensional chaotic sequences in real time on a satellite payload (a Raspberry Pi~4B) and also implement the maps on an FPGA; their output is neither a compressed stream nor authenticated~\cite{qiu2025}. Commercial JPEG encoder cores process one sample per clock at 150--166\,MHz without cryptography~\cite{logijpge,castjpeg}.

\subsection{Lightweight authenticated encryption in hardware}
Ascon, now standardized by NIST in SP~800-232~\cite{sp800232,ascon}, is small and fast in hardware: a core with one unrolled round takes about 7\,kGE, and a threshold implementation against power analysis costs about 3.1 times the area of the unprotected low-area core~\cite{gross2017}. Khairallah and Yadhunathan compare FPGA message-authentication codes built on the Ascon permutation on Artix-7~\cite{khairallah2023}. Gladis Kurian and Chen implement Ascon on an Artix-7 with replay detection through a Bloom filter in block RAM, which has no false negatives but a small false-positive rate (0.17\,\% measured)~\cite{kurian2025}.

\subsection{Traffic analysis of encrypted video}
Sizes and timing of encrypted streams reveal motion~\cite{wampler2015} and which video is played~\cite{schuster2017}; traffic shaping hides IoT activity at up to $18\times$ the bandwidth for a camera~\cite{apthorpe2019}. No encryption of content changes this, so we measure it separately (Section~\ref{sec:discussion}).

\subsection{How hardware designs in this field are evaluated}\label{sec:evalpractice}
Table~\ref{tab:evalpractice} summarizes what the hardware papers above report and how. Most stop at synthesis or place and route: Kishore et al.\ give post-synthesis resources and compute the maximum frame rate as the inverse of the coding time of one frame at the synthesized clock~\cite{kishore2015}; Ba\u{g}baba and \"Ors target a Virtex-6 ML605 kit and evaluate the encryption by PSNR and RMSE~\cite{bagbaba2015}; Khairallah and Yadhunathan report Vivado synthesis results including power~\cite{khairallah2023}; Gladis Kurian and Chen implement their design on an Arty A7-100T, show replay detection on the board with LED indicators and report on-chip power split into static and dynamic parts~\cite{kurian2025}. The most complete measurement methodology we found is the Zynq-7020 testbed of Farahmand et al.\ for cryptographic cores: an AXI DMA streams data between DDR and the core, an AXI timer measures hardware and software time to one clock cycle, every hardware output is compared with software, and the clock is raised by binary search to the highest frequency that still gives correct results~\cite{farahmand2016}. We follow the same principles: bit-exact comparison with a reference for every frame, cycle-accurate throughput, post-route timing and power with simulated switching activity, and a board protocol with the same DMA structure (Section~\ref{sec:board}).

\begin{table*}[t]
\caption{How related hardware designs were evaluated (from each source).}
\label{tab:evalpractice}
\centering\small\setlength{\tabcolsep}{4pt}
\resizebox{\textwidth}{!}{%
\begin{tabular}{lllll}
\toprule
Work & Platform & Area / timing & Throughput & Correctness, power \\
\midrule
Kishore~\cite{kishore2015} & Stratix II/IV & synthesis & computed from clock & PSNR; no power \\
Ba\u{g}baba~\cite{bagbaba2015} & Virtex-6 ML605 & synthesis & -- & PSNR, RMSE; no power \\
Khairallah~\cite{khairallah2023} & Artix-7 & Vivado synthesis & computed & Vivado power \\
Kurian~\cite{kurian2025} & Arty A7-100T board & Vivado & latency in simulation & LEDs on board; on-chip power \\
Farahmand~\cite{farahmand2016} & ZedBoard (Zynq-7020) & place and route & AXI timer, DMA & vs.\ software; max.\ frequency search \\
\midrule
Ours & PYNQ-Z2 (Zynq-7020) & post-route, OOC + system & cycle-accurate RTL + cycle counters on the board & bit-exact vs.\ model in simulation and on the board; SAIF power; board power: Vivado estimate (no meter) \\
\bottomrule
\end{tabular}}
\end{table*}

\subsection{Weaknesses of prior designs and how we address them}\label{sec:weak}
Table~\ref{tab:weak} collects, from the sources above, the weaknesses that recur in format-compliant and hardware video encryption, and the part of this work that answers each. Every answer is evaluated in the sections named in the last column; none is claimed beyond what is measured there.

\begin{table*}[t]
\caption{Recurring weaknesses of prior work (from each source) and where this work addresses them.}
\label{tab:weak}
\centering\small\setlength{\tabcolsep}{4pt}
\begin{tabular}{p{0.30\textwidth}p{0.17\textwidth}p{0.33\textwidth}p{0.11\textwidth}}
\toprule
Weakness & Seen in & This work & Where \\
\midrule
No authentication: a tampered, replayed or downgraded frame is displayed & \cite{kishore2015,bagbaba2015,qiu2025} & per-frame Ascon tag over metadata and scan, replay window, authenticated level & Sec.~\ref{sec:tag}, \ref{sec:feedback} \\
Key material that does not change per frame: one known frame breaks the rest & \cite{vandroogenbroeck2002,kishore2015,hirose2024,li2025shuffle} & keystream and E5 permutation drawn anew from (session, frame counter) & Sec.~\ref{sec:nonce}, \ref{sec:bench} \\
Field encryption leaves outlines (non-zero counts per block) & \cite{hirose2024} & E5 permutes every MCU, keyed per frame, without a frame buffer & Sec.~\ref{sec:e5}, \ref{sec:leakage} \\
Visual distortion (PSNR) used as the security measure & \cite{bagbaba2015,kishore2015,massoudi2008} & the same hand-made and learned attacks against every scheme (SVE-Bench) & Sec.~\ref{sec:bench} \\
Grey 256$\times$256 images, or a few test images & \cite{kishore2015} & colour 4:2:0 video: open films and private clips & Sec.~\ref{sec:security} \\
Throughput computed from the synthesized clock & \cite{kishore2015,khairallah2023} & cycle-accurate RTL, post-route timing, bit-exact on the board & Sec.~\ref{sec:hwresults} \\
Probabilistic replay detection (Bloom filter false positives) & \cite{kurian2025} & a sliding 64-frame window with no false positives & Sec.~\ref{sec:feedback} \\
Output not a compressed, playable stream & \cite{qiu2025} & every level is baseline JPEG (Lemma~\ref{lem:format}) & Sec.~\ref{sec:selenc} \\
\bottomrule
\end{tabular}
\end{table*}

\subsection{Gap}
We found no design that (i)~places a standardized lightweight AEAD inside a hardware JPEG encoder at the entropy-coding stage, (ii)~keeps every protection level playable while authenticating each frame together with its metadata, (iii)~hides block positions in hardware without a frame buffer, and (iv)~closes the loop from the receiver back to the encoder so that the protection level follows observed attacks. We also could not find an evaluation that scores published format-compliant schemes and new ones under the same attacks, including learned attackers, on real video.
